\section{Public Real-Code Mutation Audit}
\label{sec:real-code-audit}

The constructed mechanism cases answer RQ3 but cannot test whether the mutation shapes discussed by the paper occur in independently authored software. We therefore add a deliberately small external audit. Its purpose is not to estimate prevalence or migration effort; it asks whether selected public retained-state updates expose local shapes that the current Patch surface can represent directly, and where the present design instead requires adaptation or restructuring.

\subsection{Protocol and coding}
The corpus uses a purposive maximum-variation design fixed on 30 August 2026. Six public JavaScript/TypeScript projects are pinned to full Git commit SHAs: GitLens, Prettier for VS Code, VS Code ESLint, Obsidian Dataview, Obsidian Tasks, and Node-RED. Exactly three production retained-state mutation observations are coded per project. Tests, fixtures, generated/build-only code, benchmark-only code, and loop-local temporaries are excluded. Within each project, observations were chosen to vary mutation shape where possible rather than to maximize Patch compatibility. Consistent with software-engineering sampling guidance, this non-probability design is named explicitly and its generalization boundary is kept separate from descriptive corpus counts~\cite{baltes2022sampling}. The artifact records repository, immutable commit, path, source context, and a short auditable anchor for every observation.

Each observation is coded on four dimensions. \emph{Operation family} records the closest local Patch semantic operation (increase, decrease, set, add, remove, or clear). \emph{Local surface fit} is stricter: \emph{direct} means that the current Patch source surface can represent the local update without changing the relevant representation semantics; \emph{adapter} retains an operation-level analogue but requires a host/data-model bridge such as Set, Map, trie, dynamic keys, or host persistence; and \emph{restructure} requires a larger source transformation such as callback filtering or spread/bulk insertion. A separate \emph{Lean-fragment shape match} records only whether the local update has the same singleton numeric directional constant-magnitude form as the current Change-to-contract theorem. It does not verify the third-party program. Finally, a context code preserves whether the update is standalone, coupled across targets, host-persisted, sequential on one target, or dynamically/batch targeted.

This distinction matters. For example, JavaScript \texttt{Set.add} is locally add-like, but treating it as a direct Patch list update would erase uniqueness semantics. Likewise, a scalar counter increment can be locally direct while its enclosing method maintains a consistency relation with a separate trie. The coding retains both facts.

\subsection{Results}
Table~\ref{tab:real-code-audit} reports the deterministic summary generated from the frozen coding manifest.

\begin{table}[t]
\centering
\caption{Public real-code mutation audit. Counts are descriptive of the purposive corpus only. ``Lean shape'' means local syntactic/semantic shape alignment with the narrow singleton numeric bridge, not verification of external code.}
\label{tab:real-code-audit}
\small
\begin{tabular}{lrrrr}
\toprule
Project & Direct & Adapter & Restructure & Lean shape \\
\midrule
GitLens & 2 & 1 & 0 & 2 \\
Prettier VS Code & 0 & 3 & 0 & 0 \\
VS Code ESLint & 1 & 2 & 0 & 0 \\
Obsidian Dataview & 1 & 2 & 0 & 1 \\
Obsidian Tasks & 1 & 0 & 2 & 0 \\
Node-RED & 0 & 3 & 0 & 0 \\
\midrule
Total & 5 & 11 & 2 & 3 \\
\bottomrule
\end{tabular}
\end{table}

All 18 observations admit a local label in the coarse operation vocabulary, but this is not a portability result: only 5 are coded as direct current-surface fits, 11 require an adapter, and 2 require source restructuring. The three Lean-fragment shape matches are the GitLens repository-count increment/decrement and Dataview's revision increment. They demonstrate that the narrow numeric form is present in independently authored code, but the Lean theorem neither parses nor certifies those TypeScript programs.

The surrounding context is at least as informative as the local operation. Only 7 of 18 observations are coded standalone. The remaining 11 include four coupled multi-target contexts, two host-persisted updates, two sequential same-target rebuilds, two dynamic-target updates, and one batched dynamic-target update. Examples include GitLens maintaining a scalar repository count together with a trie, Dataview removing one file from several indices and advancing revision, VS Code ESLint writing flags through \texttt{globalState}, Obsidian Tasks rebuilding a task list by filter then spread insertion, and Node-RED selecting context-store paths dynamically. These are concrete boundaries for the present language: richer keyed collections, dynamic paths, bulk/list transforms, host-state adapters, and atomic multi-target Change would reduce the amount of restructuring needed.

\paragraph{Exploratory audit finding}
The external audit gives mixed evidence. Independently authored code contains local mutation shapes that map directly to current Patch Changes, including three observations matching the narrow numeric directional theorem shape. More observations in this deliberately varied corpus require representation/host adapters or source restructuring, and most audited contexts are not standalone. The result therefore supports external relevance of the semantic vocabulary while simultaneously identifying current language boundaries. Because the corpus is small and purposively selected, no ratio above is interpreted as ecosystem prevalence, migration cost, security improvement, or evidence that mandatory Change is superior to conventional mutation plus separate analyses.

\subsection{Research frontier exposed by the audit}
The audit also provides a concrete origin for several follow-on research questions without extending the present paper's claim set. Four observations occur inside coupled multi-target logical updates. This motivates \emph{Relational Atomic ChangeSets}, an optional future semantic unit that would stage several member Changes together and permit relations over their joint before/after state. A separate Stage-0 research branch prototypes all-or-none staging, relational checks, inverse whole-set restoration, and a Change Plan representation. It is not integrated into the evaluated stable parser, interpreter, or Change IR, and no Lean theorem currently establishes atomicity or relational invariant preservation. Those prototype properties are therefore not results of this paper.

Likewise, 11 observations require representation, dynamic-target, or host-state adapters. This motivates \emph{Certified Change Adapters}: explicit boundaries that would translate foreign retained-state operations into Patch semantic Changes while making representation assumptions and proof obligations visible. The present audit supplies examples of the need for such a mechanism, not evidence that an adapter design has already solved them.

A second isolated Stage-0 prototype derives the least capability currently expressible from an inferred Change Signature by retaining target/path, operation, and known magnitude while failing closed on unknown targets and flagging obvious excess declared authority. A research-quality \emph{least-authority inference} result still requires a formal capability ordering and proofs of sufficiency and minimality. The present manuscript therefore treats Relational ChangeSets, Certified Change Adapters, and least-authority inference as hypotheses motivated by observed pressure points, not as additional validated contributions.

\subsection{Reproducibility and coding validity}
\path{studies/real-code-mutations/corpus.json} is the frozen coding manifest and \path{results.json} is regenerated deterministically by \path{scripts/evaluate-real-code-mutations.js}. The evaluator enforces six projects, three observations per project, full commit pins, valid coding categories, and consistency of the reported counts. An optional network check resolves the immutable raw GitHub source revisions and verifies every recorded anchor; manuscript CI does not depend on network availability for those third-party sources.

The author performed the coding, with AI tools used to assist source discovery and adversarially challenge classifications. There is no independent second coder and no inter-rater reliability statistic. Exact source pins and anchors make individual judgments auditable but do not remove selection or classification bias. The study evaluates mutation shape only: it does not execute behavior-preserving translations of the six systems, measure translation effort, or compare developer outcomes.
